% Source: https://artofproblemsolving.com/community/c3h2779134_combinatorics__geometry_problem
% Stable identity: aops-topic-2779134-q01
All twelve vertices lie on a circle with centre $O$. A quadrilateral inscribed in a circle is a rectangle exactly when both of its diagonals are diameters. If it is a rectangle, each right angle is an inscribed angle standing on a diameter (Thales). Conversely, two diameters bisect each other and have equal length, so their endpoints form a rectangle.
\begin{asy}
usepackage("amsmath");
size(4cm);
defaultpen(fontsize(10pt));
pair[] V;for(int k=0;k<12;++k) V.push(dir(30*k));
draw(unitcircle,gray(0.7));
for(int k=0;k<12;++k) draw(V[k]--V[(k+1)%12]);
for(int k=0;k<6;++k) draw(V[k]--V[k+6],dotted);
draw(V[1]--V[4]--V[7]--V[10]--cycle,linewidth(1));
dot((0,0));label("$O$",(0,0),SE);
for(int k=0;k<12;++k) dot(V[k]);
\end{asy}

The twelve vertices form $6$ diametrically opposite pairs, so the number of rectangles is $\dbinom62=\boxed{15}$ (including the $3$ squares).
